\clearpage
\section*{September 21, 2026: Applying land decisions and approved earlier-floor estimates}
\textit{Agent-drafted research entry.}

The documentary review of parent boundaries (see the September 20 entry)
produced adopted ground areas. This entry records how those decisions enter
the production panels and how the earlier building floor on each reviewed
site is measured. Earlier floor matters because built FAR, the floor area of
buildings already standing on the site divided by its ground, is one of the
calibration moments behind the reweighted historical comparison. Six parents
carry an explicitly approved floor estimate; GO Broome carries a documented
zero; Wallabout carries a mapped allocation of its ground to earlier parcels.
None of these corrections changes a housing unit count, a filing date, a
refiling field or a parent membership. This entry replaces the September 21
entries on implementing the reviewed parents, the remaining three sites, the
five mergers and land diligence, and the five floor choices, together with
the GO Broome floor paragraph of the September 22 boundary batch and the
Wallabout paragraph of the September 22 Livingston follow-up.

\paragraph{How an adopted decision enters production.}
Each decision is a row of
\path{tasks/parent_opportunities_manual/output/site_lot_decisions.csv}, a
committed hand-researched table that no script edits. A row names one parent,
the constituent DOB jobs it was reviewed for
(\texttt{expected\_component\_jobs}), the named pre-filing MapPLUTO release
(\texttt{reference\_vintage}), the earlier parcels (\texttt{reference\_bbls})
and their recorded area in that release, the adopted ground
(\texttt{development\_area\_sqft}), the earlier building floor to remove
(\texttt{excluded\_building\_area\_sqft}), the measurement basis, and the
estimate flag. The site-characteristics producer
(\path{tasks/build_parent_site_characteristics}) stops if the listed jobs
differ from the parent's current constituents, if the recorded area differs
from the named release, or if the removed floor exceeds the parcels' recorded
floor. A review therefore applies only to the filing set it covered.

Residential and broad FAR always come from the earlier parcels in the named
release, never from the later successor lot. Several earlier parcels may share
one row only when they have identical residential and broad FAR; the producer
checks that equality, so no internal area split is invented. At 355 Exterior
the four contributing lots all have residential FAR 3.44 and broad FAR 6.5 and
form one row. Starhill's two old lots have residential FAR 6.02 and 3.44, so
they are separate rows and the parent's FAR is the area-weighted 5.85613.
Floor on a retained neighboring building is removed from earlier floor.

\paragraph{Earlier floor is allocated separately from ground.}
The two areas in a row are independent. Built FAR for a row is the recorded
floor of its earlier parcels, less the removed floor, divided by the adopted
ground. A partly included parcel can therefore keep all of its floor when the
building stands on the included part, or lose all of it when the building
stands on retained land. At Bedford Square the housing parent includes part of
the Sears parcel, but the full 175,875-square-foot Sears building stays on
the retained parcel; only the east auto-center's 14,946 square feet remain,
for existing FAR 0.08419 on 177,520.79 square feet. At Starhill the included
44.5-by-100-foot strip of old lot 162 (4,450 square feet) carries none of
that lot's 33,750-square-foot building, while old lot 134's 80,370 square feet
stay, giving existing FAR 1.14711 on 70,063 square feet. Motto's reference
parcels hold 83,936 square feet of floor, of which 69,357 lie on retained lot
37, leaving 14,579. At 355 Exterior the neighboring old-lot-38 building's
25,962 square feet are removed, leaving 31,850 (existing FAR 0.264816). At
Amsterdam, Casa Celina, DeKalb/LIU and Penn South the retained campus parcel
reproduces the entire earlier floor after subdivision, so all of it is
removed.

\paragraph{The six approved floor estimates.}
For six parents no document measures the partition of an earlier building
across the development boundary. Jacob approved Flatbush's estimate on
September 21 and the other five on September 21--22, each as a flagged
estimate or later administrative proxy. Their rows set
\texttt{built\_floor\_area\_estimated = TRUE}, and the canonical parent panel
carries the flag. Where the estimate is built from a proxy rather than from the
recorded floor, the removed amount is an accounting residual, not a measured
floor on retained land.

\begin{center}
\small
\begin{tabular}{llrrrr}
\hline
Parent & Sample & Units & Ground (sq ft) & Earlier floor (sq ft) & Built FAR \\
\hline
Flatbush tower & Hist. & 441 & 12,603 & 39,929.757 & 3.168274 \\
East 125th Street & Hist. & 614 & 42,540 & 50,836.28 & 1.195023 \\
St.~James/Jerome & Hist. & 102 & 17,775 & 3,650 & 0.205345 \\
Onderdonk & Post & 62 & 9,310 & 13,141.62 & 1.411560 \\
Jamaica/165th Street & Post & 591 & 39,349.50 & 40,342.49 & 1.02523 \\
Kingsbrook/Schenectady & Post & 402 & 105,382 & 226,059.00 & 2.14514 \\
\hline
\end{tabular}
\end{center}

\textit{Flatbush.} The tower's NYCIDA ground lease gives 12,603 square feet,
replacing a 1,198-square-foot parcel match. Earlier lots 23 and 24 lie wholly
inside and contribute their 2,450 and 15,640 square feet of floor in full. The
tower covers 61.3821\% of old lot 18's mapped ground, and the same share of
that lot's 35,580 square feet of floor is included, giving 39,929.757 and
removing 13,740.243. The February 2018 environmental review describes and
photographs a three-story building on lot 18, which supports a roughly uniform
height. It does not measure the partition; equal floor density across lot 18
is the assumption.

\textit{East 125th Street.} The 2021 OER work plan identifies the
42,540-square-foot site as parts of old lots 20 and 27. Each old lot's frozen
20v1 floor is multiplied by its own geometric overlap share:
$64{,}363\times0.49187+20{,}860\times0.91936=50{,}836.28$, with full-precision
shares in the calculation. Floor is assumed uniform within each old parcel.
The site was vacant by 2021, but that does not date the demolition relative
to the April 2020 filing, so zero is not adopted.

\textit{St.~James/Jerome.} The recorded deed and lease bound 17,775 square
feet. The 20v3 frozen record has 13,000 square feet of floor on the undivided
parcel. In 21v3 and 22v1, before the new building appears, the housing parcel
records 3,650 and the retained church 7,650. The adopted value is the
housing parcel's later administrative 3,650. Housing plus church is 1,700
below the frozen total, and whether that reflects a correction or a physical
change is unresolved. The removed 9,350 is the residual $13{,}000-3{,}650$.
Subtracting the church instead would have assigned 5,350.

\textit{Onderdonk.} A recorded agreement identifies developer lot 30 with
9,309.60 surveyed square feet (administrative 9,310) out of the
13,218-square-foot predecessor. Frozen 2023 floor of 18,658 is allocated by
land share: $18{,}658\times9{,}310/13{,}218=13{,}141.62$. The old building had
two stories, and later one-story DOF records cannot locate the missing second
story, so this is weaker than Flatbush's uniform-height evidence.

\textit{Jamaica/165th Street.} The 2026 recorded declaration separates the
39,349.50-square-foot residential parcel from the MTA terminal. The 2024
condominium plans locate each earlier building's shop floors, rear
mezzanines, second story and roof stair. For each old parcel, its share of
above-cellar plan area inside the boundary is applied to its frozen 2023
administrative floor. The two shops outside the drawings, old lots 98 and 99,
use their recorded 70- and 80-foot building depths. Six old parcels contribute
housing ground, with 88,365 square feet of ground and 52,593 of floor; old lots
89 and 94 (23,175 ground, 17,220 floor) are wholly outside and are counted
only in the eight-parcel audit inventory. Approximate footprints, scaled
mezzanine depths, and a similar spatial share for unmeasured cellar floor are
assumptions. Placing all three mezzanines outside or inside gives 39,678.62
to 40,735.68.

\textit{Kingsbrook/Schenectady.} A recorded agreement assigns 92,709 square
feet to Phase~I and 12,673 to its parking. HCR's corrected 2018 schedule gives
214,280 gross square feet, including basements, for the Leviton, Masin,
Blumberg and LeFrak pavilions. The June 2025 environmental assessment records
an 11,600-square-foot power plant with a cellar. A tracing of the July 2021
survey places 47.9165\% of its footprint inside Phase~I, or 5,558.31 square
feet per floor. The estimate adds that amount for the ground floor and again
for a cellar assumed to follow the footprint, plus 662.37 for half of the
labeled mezzanine bay (1,324.74 if fully covered):
$214{,}280+2\times5{,}558.31+662.37=226{,}059.00$. The half-bay is a midpoint
imputation. Omitting or filling the bay gives 225,396.63 to 226,721.37;
excluding all cellars gives 185,255.69. HCR gross area is a documented
substitute whose comparability with PLUTO's building area is unproved. The
removed 362,539.00 is the residual from the frozen 588,598 campus total, not a
measured retained-campus floor.

\paragraph{GO Broome and Wallabout.}
MapPLUTO 19v2 records 4,600 square feet of floor on GO Broome's two
development lots. City Planning dates complete demolition of the synagogue
to June 2019, before the November 2019 reference snapshot and the March 2020
filing, and the DEIS reports zero existing floor. On September 22, following
Jacob's instruction, the row removes the stale 4,600 as an explicit
source-error correction, giving zero earlier floor on the unchanged 32,395
square feet. Because the zero is documented, not imputed, the row is not
flagged as an estimate.

Wallabout's three cohort buildings occupy 39,323 square feet under the 2023
appraisal and the DOF map, separate from the earlier 251 Wallabout
development. Current successor lots are overlaid on the common pre-filing
18v2\_1 parcels, and each successor's recorded area is distributed by its
mapped overlap shares: 2,507.69 on old lot 37, 18,342.52 on old lot 41 and
18,472.79 on old lot 122. Old lot 23 contributes nothing, and the later lot
37 is a different parcel from the earlier one. All three old parcels record
zero floor. Parent ground falls from 68,805 to 39,323 square feet and
residential FAR rises from about 4.70 to 5.03. The ground total is
administrative; its split across zoning parcels is an approximate mapped
allocation. The source entries do not record a separate approval by Jacob for
this allocation.

\paragraph{Effect on the weighted comparison.}
These corrections move the reweighted historical exact-99 share by thousandths
of a percentage point and the weighted mean parent size by a few tenths of a
unit at most. The sensitivities hold outcomes, the calibration moments
and all other parents fixed, and were computed on the 554-historical,
310-post-policy sample of that date. For Flatbush on the corrected 12,603
square feet:

\begin{center}
\small
\begin{tabular}{lrrr}
\hline
Lot 18 floor included & Built FAR & Weighted exact-99 share & Weighted mean units \\
\hline
None & 1.43537 & 1.355248\% & 151.3958 \\
Adopted 61.38\% & 3.16827 & 1.354308\% & 151.4872 \\
All & 4.25851 & 1.353476\% & 151.5662 \\
\hline
\end{tabular}
\end{center}
\noindent The none-versus-all range is 0.001772 percentage points and 0.1704
units. The former 1,198-square-foot match gave 1.354249\% and 151.7519 units.

For the other five, each site's earlier floor was varied from zero to its full
frozen old total, one parent at a time (Jamaica's upper end uses the 52,593 on
its six contributing parcels). The ranges in the exact-99 share are 0.000495
percentage points (East 125th), 0.000558 (St.~James), 0.000740 (Jamaica),
0.001111 (Onderdonk) and 0.003030 (Kingsbrook). Applying all five adopted
estimates together moved the share from 1.354308\% to 1.353533\% and the mean
from 151.4872 to 151.4686 units. The September 22 batch that set GO Broome's
floor to zero and applied five boundary corrections moved them to 1.3537427\%
and 151.3559; Wallabout alone then moved the share to 1.3536772\%. These are
stress tests, not confidence intervals or joint bounds, and small effects on
these summaries neither validate a particular estimate nor establish that a
future structural model is insensitive to them.

\paragraph{Panel and sample transitions.}
Merging Bedford Square's four filings and Starhill's two reduced historical
parents from 1,818 to 1,814. Merging Godwin/Kimberly, GO Broome, Woodside,
Motto and Elara reduced historical parents to 1,811 and post-policy parents
from 909 to 907, with housing totals unchanged at 129,335 and 61,319 units.
The weighting sample fell from 557 historical and 312 post-policy parents to
554 and 310 through those five mergers alone; the floor estimates, GO Broome
and Wallabout left it unchanged. The September 22 pre-adoption source rule
later rebuilt the historical panel from the 23Q4 Housing Database and changed
the sample to 595 historical and 310 post-period parents. The
floor-specific sensitivities above have not been recomputed on that sample.
The manual table now has 63 rows for 43 parents, of which
six carry the estimate flag.

\textit{Reproduction note.} Production reads only
\path{tasks/parent_opportunities_manual/output/site_lot_decisions.csv} and the
archived MapPLUTO releases; root \texttt{make data} rebuilds the panels, and
no main task depends on an audit. The Jamaica and Kingsbrook calculations
(\path{measure_remaining_floors.R}) and the Wallabout allocation
(\path{measure_wallabout_allocation.R}) are in
\path{tasks/audits/audit_parent_site_boundaries/code/} and run under root
\texttt{make site-boundaries}. The Flatbush and five-floor sensitivity scripts
and the memos \path{five_floor_decisions.md} and
\path{flatbush_floor_review.md}, with document pages and links, were removed
from the working tree; they are preserved at tag
\texttt{pre-cleanup-2026-09-22} under
\path{tasks/audits/audit_hdb_mappluto_condo_recovery/}.
